\documentclass[11pt]{article}

\usepackage[margin=1in]{geometry}
\usepackage{amsmath,amssymb,amsthm,mathtools}
\usepackage{graphicx}
\usepackage[hidelinks]{hyperref}

\newtheorem{theorem}{Theorem}
\newtheorem{lemma}{Lemma}
\newtheorem{proposition}{Proposition}

\DeclareMathOperator{\id}{id}
\newcommand{\BH}{\mathcal B(\mathcal H)}
\newcommand{\HH}{\mathcal H}
\newcommand{\one}{\mathbf 1}
\newcommand{\ip}[2]{\langle #1,#2\rangle}

\title{Bistochastic quantum random variables over \(\mu I_{\mathcal H}\) are classical}
\author{Candidate full solution for arXiv:2005.00005}
\date{2026-06-21}

\begin{document}
\maketitle

\section*{Status}

This packet records a candidate full solution, likely valid, to the
``no example'' question raised by Plosker and Ramsey in \cite{PR} for
\(\nu=\mu I_{\HH}\).  The result covers the finite-dimensional and separable
Hilbert-space setting of the source paper.  It does not address arbitrary
positive operator-valued measures \(\nu\).

\section*{Source question}

The source paper is Sarah Plosker and Christopher Ramsey, ``Bistochastic
operators and quantum random variables'', arXiv:2005.00005 \cite{PR}.  In
Section 5, after Definition 5.1 and Theorem 5.5, the authors show that every
classical bistochastic operator on \(L^1(X,\mu)\) extends entrywise to a
bistochastic operator on \(L^1_{\HH}(X,\mu I_{\HH})\).  Immediately afterwards
they write:

\begin{center}
\includegraphics[width=0.94\linewidth]{figures/open_problem_crop.png}
\end{center}

In source-line form, the key sentence is line 1088 of the parsed source:
\[
\text{``We have no example of a bistochastic operator on }
L^1_{\HH}(X,\mu I_{\HH})\text{ that does not arise in this way.''}
\]

\section*{Definitions}

Let \((X,\Sigma,\mu)\) be a finite measure space and let \(\HH\) be finite
dimensional or separable.  Put \(\nu=\mu I_{\HH}\).  Plosker--Ramsey's
integral is entrywise in this case \cite[Section 2]{PR}: if
\(f:X\to\BH\) is a \(\nu\)-integrable quantum random variable, then
\[
        \int_X f\,d\nu = \int_X f\,d\mu
\]
in the weak operator sense, equivalently against all normal states.  They define
\(L^1_{\HH}(X,\nu)\) as the quotient of the span of positive integrable quantum
random variables by the zero seminorm, with
\[
\|f\|_1=\inf\left\{
\left\|\int_X \sum_{k=1}^4 f_k\,d\nu\right\|:
f=f_1-f_2+i(f_3-f_4),\ f_k\geq0\right\}.
\]

A linear operator \(B\) on \(L^1_{\HH}(X,\nu)\) is \emph{bistochastic} if
\[
 B\geq0,\qquad
 \int_X Bf\,d\mu=\int_X f\,d\mu\quad(f\in L^1_{\HH}(X,\nu)),
 \qquad
 B I_{\HH}= I_{\HH}.
\]
The source paper proves that a classical bistochastic operator \(S\) on
\(L^1(X,\mu)\) gives such a \(B\) by acting entrywise:
\[
       B(fA)=S(f)A,\qquad f\in L^1(X,\mu),\ A\in \BH.
\]

\section*{Proof intuition}

The first observation is that every rectangular test of an arbitrary
bistochastic \(B\) is trapped in an extremely rigid order interval.  For
measurable \(E,F\), define
\[
        \Theta_{F,E}(A)=\int_F B(\chi_E A)\,d\mu,\qquad A\in\BH .
\]
This map is positive, and because \(B\) preserves the full integral,
\[
        0\leq \Theta_{F,E}\leq \mu(E)\id_{\BH}.
\]
The positive maps in this interval are only scalar multiples of the identity
map.  In finite dimensions this follows from matrix units.  In infinite
dimensions one first handles rank-one projections and then uses domination of
vector states by pure vector states to get all of \(B(\HH)\).

Thus \(B\) is scalar on every rectangle \(\chi_E A\), and these scalars are
exactly the transition measures of a classical bistochastic operator \(S\).
The remaining difficulty is the real infinite-dimensional one: bounded/simple
operator-valued functions are not norm dense in \(L^1_{\HH}(X,\mu I_{\HH})\).
The proof avoids density.  For each \(F\), the map
\[
        T_F(f)=\int_F Bf\,d\mu
\]
is positive and dominated by \(f\mapsto\int_X f\,d\mu\).  If \(P\) is a
rank-one projection and \(Q=I-P\), then for \(f\geq0\)
\[
 \pm(PfQ+QfP)\leq \varepsilon PfP+\varepsilon^{-1}QfQ .
\]
Taking the vector state supported on \(P\), domination kills the \(QfQ\) term
and \(\varepsilon\to0\) kills the off-diagonal term.  Therefore \(T_F(f)\), in
that vector state, only sees the scalar function \(\ip{f(\cdot)\xi}{\xi}\).
Since scalar tensors are already known to be governed by \(S\), all vector
states, hence all matrix coefficients, are governed by \(S\).  This forces
\(B\) to be exactly the entrywise classical extension.

\begin{lemma}[Order interval of the identity on \(B(\HH)\)]\label{lem:order}
Let \(\HH\) be finite dimensional or separable, and let
\(\Phi:\BH\to\BH\) be a positive linear map.  If \(c\,\id-\Phi\) is positive
for some \(c\geq0\), then \(\Phi=\lambda\,\id\) for some
\(0\leq\lambda\leq c\).
\end{lemma}

\begin{proof}
The case \(c=0\) is immediate, so scale to \(c=1\).  If \(P\) is a rank-one
projection, then
\[
        0\leq \Phi(P)\leq P,
\]
hence \(\Phi(P)=\lambda_P P\) for some \(\lambda_P\in[0,1]\).

Let \(M\subseteq\HH\) be a two-dimensional subspace with orthonormal basis
\(e_1,e_2\), and let \(E_{ij}\) denote the corresponding matrix units,
viewed as operators on \(\HH\).  Applying the preceding paragraph to the
rank-one projections onto \(\mathbb C(e_1\pm e_2)\) gives
\[
 \Phi(E_{11}+E_{22}\pm E_{12}\pm E_{21})
 =\alpha_\pm(E_{11}+E_{22}\pm E_{12}\pm E_{21})
\]
after multiplying by the harmless common factor \(2\).  Adding the plus and
minus equations and comparing with
\(\Phi(E_{11})+\Phi(E_{22})=\lambda_{E_{11}}E_{11}+\lambda_{E_{22}}E_{22}\)
gives \(\alpha_+=\alpha_-\) and
\(\lambda_{E_{11}}=\lambda_{E_{22}}\).  The same argument with
\(e_1\pm i e_2\) also gives the corresponding identities for the skew
matrix-unit combination.  Consequently \(\Phi\) is a scalar multiple of the
identity on every \(2\times2\) corner.

It follows that all rank-one projections have one common scalar, say
\(\lambda\): any two rank-one projections lie in a common subspace of
dimension at most two.  Now fix a unit vector \(\xi\), let \(P_\xi\) be the
projection onto \(\mathbb C\xi\), and define the positive functional
\[
        \varphi_\xi(A)=\ip{\Phi(A)\xi}{\xi}.
\]
Since \(0\leq\Phi\leq\id\), one has
\(0\leq\varphi_\xi\leq \omega_\xi\), where
\(\omega_\xi(A)=\ip{A\xi}{\xi}\).  A positive functional dominated by the
vector state \(\omega_\xi\) is a scalar multiple of \(\omega_\xi\): indeed,
if \(P=P_\xi\), then \(\varphi_\xi(I-P)=0\), and the Cauchy--Schwarz
inequality for positive functionals gives
\(\varphi_\xi((I-P)A)=\varphi_\xi(A(I-P))=0\); hence
\[
        \varphi_\xi(A)=\varphi_\xi(PAP)=\varphi_\xi(P)\omega_\xi(A).
\]
But \(\varphi_\xi(P)=\ip{\Phi(P)\xi}{\xi}=\lambda\).  Thus
\[
        \ip{\Phi(A)\xi}{\xi}=\lambda\ip{A\xi}{\xi}
\]
for every \(A\in\BH\) and every unit vector \(\xi\).  Polarization gives
\(\Phi(A)=\lambda A\) for all \(A\).
\end{proof}

\begin{lemma}[Rectangular scalarization]\label{lem:rectangles}
Let \(B\) be a bistochastic operator on \(L^1_{\HH}(X,\mu I_{\HH})\).  For
each \(F\in\Sigma\) there is a function \(h_F\in L^\infty(X,\mu)\) with
\(0\leq h_F\leq1\) such that
\[
        \int_F B(gA)\,d\mu=\left(\int_X h_F g\,d\mu\right)A
        \tag{1}\label{eq:TFscalar}
\]
for every \(g\in L^1(X,\mu)\) and every \(A\in\BH\).
\end{lemma}

\begin{proof}
For \(E,F\in\Sigma\), set
\[
        \Theta_{F,E}(A)=\int_F B(\chi_E A)\,d\mu,\qquad A\in\BH .
\]
If \(A\geq0\), then \(\chi_EA\geq0\), so \(\Theta_{F,E}(A)\geq0\).  Moreover,
\[
        \Theta_{X,E}(A)
        =\int_X B(\chi_EA)\,d\mu
        =\int_X \chi_EA\,d\mu
        =\mu(E)A,
\]
and
\[
        \mu(E)A-\Theta_{F,E}(A)=\Theta_{X\setminus F,E}(A)\geq0
        \qquad(A\geq0).
\]
By Lemma~\ref{lem:order}, \(\Theta_{F,E}=a(E,F)\id_{\BH}\) for a scalar
\(a(E,F)\in[0,\mu(E)]\).

Fix \(F\).  The set function \(E\mapsto a(E,F)\) is a finite positive measure:
countable additivity follows by applying the normal vector state
\(\omega_\xi\) to
\(\int_F B(\chi_E I_{\HH})\,d\mu=a(E,F)I_{\HH}\).  Also
\(a(E,F)\leq\mu(E)\), so \(a(\cdot,F)\ll\mu\) and its Radon--Nikodym density
\[
        h_F=\frac{d\,a(\cdot,F)}{d\mu}
\]
satisfies \(0\leq h_F\leq1\).  Therefore \eqref{eq:TFscalar} holds first for
indicator functions \(g=\chi_E\), hence by linearity for scalar simple
functions \(g\).

Finally let \(g\in L^1(X,\mu)\).  Choose scalar simple \(g_n\to g\) in
\(L^1\).  Plosker--Ramsey prove that multiplication by a fixed
\(A\in\BH\) sends \(L^1_{\HH}(X,\mu I_{\HH})\) boundedly to itself
\cite[Proposition 4.3]{PR}; hence \(g_nA\to gA\) in \(L^1_{\HH}\).  The map
\[
        T_F(u)=\int_F Bu\,d\mu
\]
is bounded, for instance with \(\|T_Fu\|\leq2\|u\|_1\): for \(u\geq0\) this follows from
\(0\leq T_Fu\leq\int_Xu\,d\mu\), and the general case follows by taking the
infimum over positive four-term decompositions in the definition of
\(\|\cdot\|_1\).  Passing to the limit gives \eqref{eq:TFscalar}.
\end{proof}

\begin{lemma}[Dominated localization]\label{lem:localization}
With \(B,T_F\), and \(h_F\) as in Lemma~\ref{lem:rectangles}, one has
\[
        \int_F Bf\,d\mu=\int_X h_F f\,d\mu
        \tag{2}\label{eq:localization}
\]
for every \(f\in L^1_{\HH}(X,\mu I_{\HH})\).
\end{lemma}

\begin{proof}
It is enough by linearity to prove the claim for \(f\geq0\).  Fix such an
\(f\), a unit vector \(\xi\in\HH\), and let \(P\) be the rank-one projection
onto \(\mathbb C\xi\), \(Q=I-P\).  Write
\[
        a(x)=\ip{f(x)\xi}{\xi}.
\]
The constant-operator multiplier result \cite[Proposition 4.3]{PR} ensures
that \(PfP\), \(QfQ\), and \(PfQ+QfP\) all belong to \(L^1_{\HH}\).
Then \(PfP=aP\), and Lemma~\ref{lem:rectangles} gives
\[
        \ip{T_F(PfP)\xi}{\xi}
        = \ip{T_F(aP)\xi}{\xi}
        = \int_X h_F a\,d\mu.
        \tag{3}\label{eq:pfp}
\]

Since \(QfQ\geq0\), positivity and integral domination give
\[
        0\leq T_F(QfQ)\leq \int_X QfQ\,d\mu.
\]
Taking the vector state at \(\xi\) yields
\[
        0\leq \ip{T_F(QfQ)\xi}{\xi}
        \leq \int_X \ip{Qf(x)Q\xi}{\xi}\,d\mu=0.
        \tag{4}\label{eq:qfqzero}
\]

For every \(\varepsilon>0\), the pointwise block-operator inequalities
\[
        \pm(PfQ+QfP)\leq \varepsilon PfP+\varepsilon^{-1}QfQ
        \tag{5}\label{eq:blockineq}
\]
follow from positivity of
\[
        (\sqrt{\varepsilon}P\pm \varepsilon^{-1/2}Q)
        f
        (\sqrt{\varepsilon}P\pm \varepsilon^{-1/2}Q).
\]
Applying the positive map \(T_F\), taking the vector state at \(\xi\), and
using \eqref{eq:qfqzero}, we get
\[
 \left|\ip{T_F(PfQ+QfP)\xi}{\xi}\right|
 \leq
 \varepsilon\,\ip{T_F(PfP)\xi}{\xi}.
\]
Letting \(\varepsilon\downarrow0\) gives
\[
        \ip{T_F(PfQ+QfP)\xi}{\xi}=0.
        \tag{6}\label{eq:offdiagzero}
\]

Combining \eqref{eq:pfp}, \eqref{eq:qfqzero}, and
\eqref{eq:offdiagzero}, we obtain
\[
        \ip{T_F(f)\xi}{\xi}
        =
        \int_X h_F\ip{f(x)\xi}{\xi}\,d\mu
        =
        \ip{\left(\int_X h_F f\,d\mu\right)\xi}{\xi}.
\]
Here \(h_Ff\in L^1_{\HH}\) because scalar \(L^\infty\)-multipliers are bounded
on \(L^1_{\HH}\) in the case \(\nu=\mu I_{\HH}\) \cite[Lemma 4.1]{PR}.  Since
the displayed equality holds for every unit vector \(\xi\), the two positive
operators in \eqref{eq:localization} have the same quadratic forms.  Hence
\eqref{eq:localization} holds for positive \(f\), and therefore for all
\(f\) by decomposing into four positive parts.
\end{proof}

\begin{theorem}\label{thm:full}
Let \((X,\Sigma,\mu)\) be a finite measure space and let \(\HH\) be finite
dimensional or separable.  If \(B\) is a bistochastic operator on
\(L^1_{\HH}(X,\mu I_{\HH})\) in the sense of \cite{PR}, then there is a unique
classical bistochastic operator \(S\) on \(L^1(X,\mu)\) such that
\[
        B(fA)=S(f)A
        \qquad
        (f\in L^1(X,\mu),\ A\in\BH).
        \tag{7}\label{eq:tensor}
\]
Equivalently, for every \(F\in L^1_{\HH}(X,\mu I_{\HH})\), \(B(F)\) is obtained
by applying \(S\) to all scalar matrix coefficients of \(F\).  Consequently,
every bistochastic operator on \(L^1_{\HH}(X,\mu I_{\HH})\) is one of the
entrywise classical extensions constructed in \cite[Theorem 5.5]{PR}.
\end{theorem}

\begin{proof}
By Lemma~\ref{lem:rectangles}, \(B\) maps scalar simple tensors \(gI_{\HH}\)
to scalar-valued functions times \(I_{\HH}\).  The scalar subspace
\(L^1(X,\mu)I_{\HH}\) is norm closed in \(L^1_{\HH}(X,\mu I_{\HH})\)
\cite[Section 5, before Theorem 5.5]{PR}; since \(B\) is \(L^1\)-contractive
\cite[Proposition 5.4]{PR}, it follows that for every \(g\in L^1(X,\mu)\)
there is a unique scalar \(Sg\in L^1(X,\mu)\) such that
\[
        B(gI_{\HH})=(Sg)I_{\HH}.
\]
This defines a positive linear operator \(S:L^1(X,\mu)\to L^1(X,\mu)\).  Since
\(BI_{\HH}=I_{\HH}\), \(S1=1\).  Since \(B\) preserves integrals,
\[
        \int_X Sg\,d\mu=\int_X g\,d\mu\qquad(g\in L^1(X,\mu)).
\]
Thus \(S\) is a classical bistochastic operator.

For this \(S\), Lemma~\ref{lem:rectangles} says that
\[
        \int_F Sg\,d\mu=\int_X h_Fg\,d\mu
        \qquad(F\in\Sigma,\ g\in L^1(X,\mu)).
        \tag{8}\label{eq:adjoint}
\]
Indeed, apply \eqref{eq:TFscalar} with \(A=I_{\HH}\) and compare with
\(\int_F B(gI_{\HH})\,d\mu\).

Now let \(f\in L^1_{\HH}(X,\mu I_{\HH})\).  Lemma~\ref{lem:localization} and
\eqref{eq:adjoint} give, for every \(F\in\Sigma\) and every
\(\xi,\eta\in\HH\),
\[
\begin{aligned}
 \ip{\left(\int_F Bf\,d\mu\right)\xi}{\eta}
 &= \int_X h_F(x)\ip{f(x)\xi}{\eta}\,d\mu(x) \\
 &= \int_F S\big(\ip{f(\cdot)\xi}{\eta}\big)\,d\mu .
\end{aligned}
        \tag{9}\label{eq:coefficients}
\]
The scalar coefficient \(\ip{f(\cdot)\xi}{\eta}\) is in \(L^1(X,\mu)\) by
Plosker--Ramsey's state-topology estimate \cite[Lemma 3.3]{PR} and
polarization.

Let \(C\) denote the entrywise extension of \(S\) constructed in
\cite[Theorem 5.5]{PR}.  Equation \eqref{eq:coefficients} says that every
scalar matrix coefficient of \(\int_F Bf\,d\mu\) equals the corresponding
coefficient of \(\int_F Cf\,d\mu\).  Hence
\[
        \int_F Bf\,d\mu=\int_F Cf\,d\mu\qquad(F\in\Sigma).
\]
Taking scalar coefficients over a countable dense subset of the separable
Hilbert space and using equality of scalar integrals over all \(F\), we get
\(Bf=Cf\) almost everywhere, hence as elements of \(L^1_{\HH}(X,\mu I_{\HH})\).
This proves that \(B\) is the entrywise extension of \(S\).

In particular, for scalar \(f\in L^1(X,\mu)\) and \(A\in\BH\), this gives
\eqref{eq:tensor}.  Uniqueness of \(S\) is immediate from
\(B(fI_{\HH})=(Sf)I_{\HH}\).
\end{proof}

\section*{Verification and novelty check}

No computational verification is required.  The proof reduces the problem to
three checkable ingredients: the order interval
\([0,c\id]\) in positive maps on \(B(\HH)\), the rectangular scalarization
argument, and the dominated-localization lemma that bypasses the failure of
\(L^\infty_{\HH}\)-norm-density in infinite dimension.

Local run indexes were searched for arXiv id \(2005.00005\), the title
keywords, ``bistochastic'', ``quantum random variables'', ``Plosker'',
``Ramsey'', and the exact source phrase ``We have no example of a bistochastic
operator''.  The only duplicate was the earlier finite-dimensional partial
packet, now kept in this packet's history folder.

A bounded web search on 2026-06-21 for the exact source sentence and close
phrases around \(L^1_{\HH}(X,\mu I_{\HH})\), ``bistochastic operators'', and
``quantum random variables'' found the source arXiv record \cite{PR} and no
later paper explicitly answering the question.

Human review should focus on Lemma~\ref{lem:localization}, especially the use
of the block inequalities \eqref{eq:blockineq}, and on the final passage from
all matrix-coefficient integrals to equality in \(L^1_{\HH}\).

\section*{Limitations}

The theorem resolves the source-paper question for the classical POVM
\(\nu=\mu I_{\HH}\), which is exactly the setting of the ``no example''
sentence.  It does not classify bistochastic operators for arbitrary POVMs
\(\nu\), and it does not address majorization questions beyond this
classification.

\begin{thebibliography}{9}
\bibitem{PR}
S. Plosker and C. Ramsey,
\emph{Bistochastic operators and quantum random variables},
arXiv:2005.00005, 2020.
\end{thebibliography}

\end{document}
